% Abstract — The Venus Ring of Voynich f57v
% Approximately 500 words. Long-form preprint allows extended abstract.

The \vms\ (Beinecke MS~408, radiocarbon-dated 1404--1438~CE) has resisted decipherment for over a century. Two companion studies by the present author have established (i)~that nineteen substitution-based decipherment models against five Semitic languages all converge on the same structural ceiling without producing readable text~\citep{jimenez2026negative}, and (ii)~that the manuscript's text exhibits a two-tier positional encoding consistent with an encyclopaedic code rather than a natural language in disguise~\citep{jimenez2026twotier}. Both studies characterised the manuscript's global structure but neither identified a specific historical or textual tradition to which the manuscript belongs.

This third study reverses the analytical strategy: rather than analysing global structure, it identifies and quantitatively characterises a single \emph{local anomaly} in the manuscript and traces it to a documented historical tradition.

We report that folio f57v---a cosmological circular diagram showing four nymphs around a central solar symbol~$\odot$---contains a \emph{private subalphabet} of five glyphs (X2, Xd, Xf, Pc, Ea) that appear nowhere else in the entire 37,087-word STA1 corpus, plus three strongly enriched glyphs (Eb, Ja, X1; enrichment factors 60$\times$--306$\times$ over the corpus baseline). The probability of this concentration occurring by chance is approximately~$10^{-26}$. The subalphabet is concentrated in the four concentric text rings of the folio (loci tagged \texttt{@Cc}/\texttt{+Cc} in the IVTFF format), most densely in line f57v.3---the famous ``second ring of 17 characters $\times$ 4 cycles'' identified intuitively by voynichological scholarship for decades. Folio f57v is the \emph{only} folio in the 224-folio corpus with this structural signature; a robustness test against the entire cosmological/astronomical iconographic family of the manuscript (the 14 foldout panels of ff.~67--70, totalling $7134$ comparison atoms, including all the anthropomorphic sun and moon diagrams) finds \emph{zero} occurrences of the seven strongest subalphabet glyphs and only two stray tokens of the weakest one (Ja), confirming that the phenomenon is not a feature of the cosmological page genre but a localized inscription specific to f57v.

Extending the positional categorial encoding system previously validated by the author's work on the zodiac block to the labels of f57v's central nymphs (\texttt{@L0}) and to the radial labels of f70r1's sun rays (\texttt{@Ri}), we further establish that the manuscript's general labeling system follows a unified canonical morphology across cosmological and zodiacal folios. The Voynich therefore contains \emph{two textually distinct phenomena} that coexist without contaminating each other: a general categorial labeling system (Phenomenon~A) covering most of the manuscript, and a single localized enumerative inscription (Phenomenon~B) on the Venus ring of f57v.

Direct verification against the photographic facsimile of Cracow Jagiellonian Library MS~793---the only extant fully illustrated copy of \emph{\msc}, dated 1471--1474 with a reconstructed Bohemian prototype c.~1390--1410 from the workshop of Wac\l aw~IV---yields three iconographic correspondences with f57v: (i)~the central solar disc of f57v is iconographically identical to the disc held by the right-hand Sun figure on MS~793 f.~191r (folio~382 of the facsimile); (ii)~the geometrized signs that \citet{sniezynska2009} describes as ``transformations of the Arabic or Semitic alphabet'' are visible on the shields of both Sun figures on f.~191r; (iii)~a polar unwrapping of the f57v second ring reveals atomic characters of the same abstract graphic genre as the Picatrix shield sigils, suggesting that the Voynich ring contains the alphabet of which the compound Picatrix sigils could be constructed.

We propose that the \vms\ is an \emph{independent witness} of the medieval astrological-talismanic tradition that \emph{\msc} codifies, not a textual derivation. The Voynich preserves the \emph{karacteres} alphabet \emph{more completely} than MS~793 itself, where the scribe of 1471--1474 had already suppressed most of the karacteres prescribed by the original Picatrix text. This is consistent with the Voynich's earlier C14 dating and with the documented Toledo$\rightarrow$Prague$\rightarrow$Cracow transmission chain of Picatrix tradition active in Cracow from at least 1410. The work also supports a positive identification of folio f57v as a \emph{solar-year diagram} with the four inner planets (Sun, Venus, Mercury, Moon) as concentric rings and the four cardinal seasons as central nymphs.
